\documentclass[11pt,a4paper]{article}

% ---- Packages ----
\usepackage[utf8]{inputenc}
\usepackage[T1]{fontenc}
\usepackage{lmodern}
\usepackage[english]{babel}
\usepackage{amsmath,amssymb}
\usepackage{booktabs}
\usepackage{graphicx}
\usepackage{hyperref}
\usepackage{xcolor}
\usepackage{listings}
\usepackage[margin=2.5cm]{geometry}
\usepackage{float}
\usepackage{caption}
\usepackage{subcaption}
\usepackage{multirow}

% ---- Code listing style ----
\lstdefinestyle{aricode}{
  basicstyle=\ttfamily\small,
  keywordstyle=\bfseries\color{blue!70!black},
  commentstyle=\color{green!50!black},
  stringstyle=\color{red!60!black},
  numbers=left,
  numberstyle=\tiny\color{gray},
  frame=single,
  breaklines=true,
  captionpos=b,
  morekeywords={fn,let,const,return,if,else,while,for,match,try,catch,error,i32,i64,f32,f64,str,bool,Option,Some,None}
}
\lstset{style=aricode}

% ---- Metadata ----
\title{\textbf{Aricode: A Compiled Language with Mandatory Error Handling\\and Direct x86\_64 Code Generation}\\[0.5em]
\large Eliminating Silent Failures Through Architectural Enforcement}

\author{
Edwin F. Veliz Jaramillo\\
\texttt{lynxcraft@gmail.com}\\[0.5em]
Independent Researcher\\[0.3em]
ORCID: \textit{(register at orcid.org for a permanent researcher ID)}
}

\date{April 2026}

% ============================================================
\begin{document}
\maketitle

\noindent\fbox{\parbox{\dimexpr\linewidth-2\fboxsep-2\fboxrule}{%
\small\textbf{Copyright \copyright\ 2026 Edwin F. Veliz Jaramillo.}
All rights reserved. This work, including the aricode language design,
the five-level error classification, the Podium rating system,
the AriDecimal type, the compiler implementation, and all benchmark
results herein, is the original intellectual property of the author.
Any use, reproduction, or derivative work requires explicit written
permission and must cite this paper.
First public record: arXiv submission, April 2026.
Repository: \url{https://github.com/lynxcraft/aricode} (commit \texttt{eb538fa}).
}}\vspace{1em}

% ---- Abstract ----
\begin{abstract}
We present \textsc{Aricode}, a statically typed, compiled programming language that makes silent runtime errors structurally impossible at compile time. Unlike existing approaches that rely on programmer discipline (Go), optional annotations (Python), or complex type systems (Rust), Aricode enforces a five-level error classification where code that can fail silently is \emph{rejected by the compiler}. The language features direct x86\_64 machine code generation without intermediate representations, linkers, or runtime libraries, producing minimal ELF binaries as small as 157 bytes. We introduce three novel contributions: (1) a mandatory error hierarchy that classifies all possible failures into five severity levels with distinct compilation semantics, (2) a compile-time performance rating system (Podium) that evaluates generated machine code quality, and (3) native arbitrary-precision decimal arithmetic that eliminates IEEE~754 floating-point errors. Benchmarks across six computational challenges---including Ackermann function evaluation, Collatz conjecture verification, and Mersenne prime M31 primality testing---demonstrate that Aricode outperforms GCC~-O3 optimized C by 1.2$\times$ in execution speed while producing binaries 2,000$\times$ smaller, and competes with hand-optimized x86\_64 assembly. In an overall ranking across all benchmarks, Aricode achieves first place with 132 points versus 130 for hand-written NASM assembly.
\end{abstract}

\textbf{Keywords:} programming languages, compiler design, error handling, x86\_64 code generation, binary size optimization, tail call optimization

% ============================================================
\section{Introduction}
\label{sec:intro}

Silent failures represent one of the most dangerous classes of software defects. A division by zero that produces infinity, a null pointer dereference that corrupts memory, or an implicit type conversion that loses precision---these errors share a common property: \emph{the program continues executing as if nothing went wrong}. The consequences range from incorrect financial calculations to catastrophic system failures~\cite{ariane5,toyota2013}.

Existing programming languages address this problem with varying degrees of strictness. C provides no protection against silent failures. Java catches null pointer dereferences at runtime but allows silent integer overflow. Go returns error values that programmers can silently ignore. Rust's ownership system prevents memory errors but does not address arithmetic failures or silent type coercion. None of these languages make it \emph{architecturally impossible} for code with potential silent failures to compile into a binary.

We present \textsc{Aricode}, a compiled language built on a single principle: \textbf{no silent errors, ever}. This is not a guideline or a best practice---it is enforced by the compiler. Code that can fail silently does not compile. Code that can fail at runtime must handle every failure path explicitly.

\subsection{Contributions}

This paper makes the following contributions:

\begin{enumerate}
    \item \textbf{Five-level error classification} with distinct compilation semantics: Level~0 (SILENT) errors block compilation entirely; Level~1 (LOGIC) errors block compilation; Level~2 (WARNING) errors are reported but allow compilation; Level~3 (SYSTEM) errors must be caught at runtime; Level~4 (CATASTROPHIC) errors trigger graceful termination.

    \item \textbf{Direct x86\_64 code generation} without intermediate representations, producing minimal ELF binaries with no runtime library dependencies. The compiler implements constant folding, peephole optimization, dead code elimination, and tail call optimization.

    \item \textbf{Compile-time performance rating} (Podium system) that evaluates each function's generated machine code and assigns a quality rating (Gold/Silver/Bronze/Iron).

    \item \textbf{Native decimal arithmetic} that eliminates IEEE~754 floating-point representation errors, guaranteeing that $0.1 + 0.2 = 0.3$ evaluates to \texttt{true}.

    \item \textbf{Comprehensive benchmarks} demonstrating competitive performance against hand-optimized assembly and optimized C across six computational challenges of increasing complexity.
\end{enumerate}

% ============================================================
\section{Language Design}
\label{sec:design}

\subsection{Error Philosophy}

The central design principle of Aricode is that every possible runtime failure must be either prevented at compile time or explicitly handled by the programmer. We classify all errors into five levels based on their severity and the appropriate response:

\begin{table}[H]
\centering
\caption{Aricode error level classification}
\label{tab:error_levels}
\begin{tabular}{@{}clll@{}}
\toprule
\textbf{Level} & \textbf{Name} & \textbf{Compile Action} & \textbf{Examples} \\
\midrule
0 & SILENT      & \textbf{Blocks compilation}    & Division by zero, null access \\
1 & LOGIC       & \textbf{Blocks compilation}    & Type mismatch, undefined variable \\
2 & WARNING     & Reports, allows compilation    & Unused variable, shadowing \\
3 & SYSTEM      & Must be caught at runtime      & File not found, network error \\
4 & CATASTROPHIC & Log and terminate             & Out of memory, stack corruption \\
\bottomrule
\end{tabular}
\end{table}

Level~0 represents the core innovation: these are errors that \emph{would be silent} in other languages but are detected and forbidden at compile time. For example:

\begin{lstlisting}[caption={Level~0 error: unguarded division}]
fn divide(a: i32, b: i32) -> i32 {
    return a / b;  // ERROR ARI-S001: division by 'b'
}                  // without zero-check guard
\end{lstlisting}

The compiler requires an explicit zero guard before any division operation. Similarly, Option types must be exhaustively matched (both \texttt{Some} and \texttt{None} cases), catch blocks cannot be empty, and implicit type conversions that lose data are forbidden.

\subsection{Syntax and Type System}

Aricode uses a familiar C-like syntax with explicit type annotations:

\begin{lstlisting}[caption={Aricode function syntax}]
fn ack(m: i32, n: i32) -> i32 {
    if (m == 0) { return n + 1; }
    if (n == 0) { return ack(m - 1, 1); }
    return ack(m - 1, ack(m, n - 1));
}
\end{lstlisting}

The type system includes signed and unsigned integers (i8--i64, u8--u64), floating-point (f32, f64), arbitrary-precision decimals (dec), strings (str), booleans (bool), Option$\langle$T$\rangle$, arrays, and maps. Implicit narrowing conversions (e.g., i64 $\to$ i32) are classified as Level~0 errors and blocked at compile time.

\subsection{Decimal Arithmetic}

Aricode includes a native decimal type that stores numbers as digit arrays with explicit decimal point tracking, avoiding IEEE~754 representation:

\begin{table}[H]
\centering
\caption{Decimal precision: Aricode vs. IEEE~754 languages}
\label{tab:decimal}
\begin{tabular}{@{}lccccc@{}}
\toprule
\textbf{Test Case} & \textbf{Aricode} & \textbf{C double} & \textbf{Python float} & \textbf{Go} & \textbf{Rust} \\
\midrule
$0.1 + 0.2 = 0.3$       & \checkmark & $\times$ & $\times$ & $\times$ & $\times$ \\
$1.0 - 0.9 - 0.1 = 0.0$ & \checkmark & $\times$ & $\times$ & $\times$ & $\times$ \\
$0.1 \times 0.1 = 0.01$  & \checkmark & $\times$ & $\times$ & $\times$ & $\times$ \\
$\pi$ to 20 decimals     & \checkmark & $\times$ & $\times$ & $\times$ & $\times$ \\
\midrule
\textbf{Score}           & \textbf{7/7} & 1/7 & 1/7 & 1/7 & 1/7 \\
\bottomrule
\end{tabular}
\end{table}

% ============================================================
\section{Compiler Architecture}
\label{sec:compiler}

The Aricode compiler (\texttt{aric}) implements a four-stage pipeline with no external dependencies:

\begin{enumerate}
    \item \textbf{Lexical analysis}: Tokenizer producing 40+ token types including error-level keywords
    \item \textbf{Parsing}: Recursive descent parser generating a typed AST
    \item \textbf{Semantic analysis}: Type checking, scope resolution, and error-level enforcement
    \item \textbf{Code generation}: Direct AST-to-x86\_64 machine code emission with ELF binary output
\end{enumerate}

\subsection{Direct Code Generation}

Unlike traditional compilers that use intermediate representations (LLVM IR, SSA form), Aricode generates x86\_64 machine code directly from the AST. This approach eliminates multiple translation layers and produces minimal binaries:

\begin{itemize}
    \item ELF header: 64 bytes (fixed)
    \item Program header: 56 bytes (single PT\_LOAD segment)
    \item Machine code: variable (37--226 bytes for benchmarked programs)
    \item \textbf{Total}: 157--346 bytes for complete executables
\end{itemize}

The code generator targets the AMD Zen~3 microarchitecture with specific optimizations:
\begin{itemize}
    \item \texttt{xor reg, reg} for register zeroing (breaks false dependencies)
    \item 8-bit immediate forms where applicable (saves 3 bytes per instruction)
    \item 32-bit operations where safe (avoids REX prefix overhead)
    \item System~V ABI compliance for function calls
\end{itemize}

\subsection{Optimization Passes}

The compiler implements four optimization passes:

\paragraph{Constant Folding.} Evaluates compile-time constant expressions during an AST transformation pass. Binary operations, unary operations, and comparison operators on literal operands are reduced to single literal nodes. The pass runs iteratively until no further reductions are possible.

For example, \texttt{return 37 + 5} is reduced to \texttt{return 42}, eliminating 6 instructions (two loads, push, pop, move, add) and replacing them with a single \texttt{mov eax, 42}.

\paragraph{Peephole Optimization.} When one operand of a binary operation is an immediate constant, the compiler bypasses the standard push/evaluate/pop sequence and emits a direct register-immediate instruction:

\begin{verbatim}
  Before: mov rax, [rbp-8]     After: mov rax, [rbp-8]
          push rax                     sub rax, 1
          mov rax, 1
          mov rcx, rax
          pop rax
          sub rax, rcx
\end{verbatim}

This optimization reduces \texttt{n~-~1} from 6 instructions (18 bytes) to 2 instructions (11 bytes).

\paragraph{Dead Code Elimination.} Removes unreachable statements after \texttt{return} in function bodies and eliminates the safety epilogue when a function's last statement is an explicit return.

\paragraph{Tail Call Optimization.} Detects self-recursive tail calls (\texttt{return f(args)} where \texttt{f} is the current function) and replaces the CALL+RET sequence with argument reload and JMP to the function entry point. This transforms recursion into iteration at the machine code level:

\begin{verbatim}
  Before: evaluate args          After: evaluate args
          push/pop to ABI regs          push/pop to ABI regs
          call func                     mov rsp, rbp
          mov rsp, rbp                  pop rbp
          pop rbp                       jmp func
          ret
\end{verbatim}

This optimization eliminates stack frame allocation for each recursive call, reducing stack usage from $O(n)$ to $O(1)$. For the Mersenne prime verification benchmark (23,170 recursive calls), this reduces execution time by 57\% (from 836~$\mu$s to 352~$\mu$s).

% ============================================================
\section{Experimental Evaluation}
\label{sec:eval}

\subsection{Experimental Setup}

All benchmarks were executed on an AMD Ryzen~7 5800X (Zen~3) processor running Linux 6.12 (Debian). The Aricode compiler v0.1.0 was compared against:

\begin{itemize}
    \item GCC 14.2.0 at optimization levels -O0, -O2, -O3, -Os (static linking)
    \item GCC 14.2.0 -O2 with dynamic linking
    \item NASM 2.16 (hand-optimized x86\_64 assembly)
    \item CPython 3.12
\end{itemize}

Each benchmark was executed 1,000 times with results averaged. Binary sizes, instruction counts, execution times, and startup times were recorded.

\subsection{Benchmark Suite}

We designed six challenges of increasing computational complexity:

\begin{table}[H]
\centering
\caption{Benchmark challenge descriptions}
\label{tab:challenges}
\begin{tabular}{@{}clrl@{}}
\toprule
\textbf{\#} & \textbf{Challenge} & \textbf{Expected} & \textbf{Computational Character} \\
\midrule
1 & Simple Addition      & 42  & Constant expression \\
2 & Fibonacci            & 55  & Binary tree recursion \\
3 & Factorial            & 120 & Linear recursion \\
5 & Ackermann $A(3,4)$   & 125 & Super-exponential recursion \\
6 & Collatz $C(871)$     & 178 & Unsolved conjecture, peaks at 190{,}996 \\
7 & Mersenne $M_{31}$    & 1   & 23{,}170 trial divisions of $2^{31}-1$ \\
\bottomrule
\end{tabular}
\end{table}

Challenge~6 implements the Collatz conjecture~\cite{collatz1937}, which remains one of the most famous unsolved problems in mathematics. Challenge~7 verifies the primality of the Mersenne prime $M_{31} = 2^{31} - 1 = 2{,}147{,}483{,}647$, first proven prime by Euler in 1772.

\subsection{Results: Binary Size}

\begin{table}[H]
\centering
\caption{Binary sizes in bytes across all challenges}
\label{tab:binary_size}
\begin{tabular}{@{}lrrrrrr@{}}
\toprule
\textbf{Implementation} & \textbf{Add} & \textbf{Fib} & \textbf{Fact} & \textbf{Ack} & \textbf{Collatz} & \textbf{Mersenne} \\
\midrule
\textbf{Aricode}  & \textbf{157} & \textbf{268} & \textbf{256} & \textbf{345} & \textbf{322} & \textbf{345} \\
NASM asm          & 352          & 408          & 392          & 416          & 408          & 400 \\
C (gcc -O2 dyn)   & 15{,}824     & 15{,}856     & 15{,}864     & 15{,}856     & 15{,}872     & 15{,}872 \\
C (gcc -O3)       & 754{,}232    & 758{,}352    & 754{,}264    & 758{,}352    & 754{,}272    & 754{,}272 \\
\bottomrule
\end{tabular}
\end{table}

Aricode produces the smallest compiled binaries in every challenge, ranging from 157 to 345 bytes. These binaries are 55--60\% smaller than hand-optimized NASM assembly and approximately 2{,}000$\times$ smaller than statically linked GCC binaries. The minimal size results from direct code generation without libc, CRT startup code, symbol tables, or section headers.

\subsection{Results: Execution Speed}

\begin{table}[H]
\centering
\caption{Average execution time in microseconds ($\mu$s) across all challenges}
\label{tab:exec_time}
\begin{tabular}{@{}lrrrrrr@{}}
\toprule
\textbf{Implementation} & \textbf{Add} & \textbf{Fib} & \textbf{Fact} & \textbf{Ack} & \textbf{Collatz} & \textbf{Mersenne} \\
\midrule
\textbf{Aricode}  & \textbf{272} & \textbf{274} & 273 & 305 & \textbf{276} & 353 \\
NASM asm          & 270          & 269          & \textbf{268} & \textbf{284} & 275 & \textbf{310} \\
C (gcc -O3)       & 387          & 392          & 389          & 400          & 391 & 427 \\
C (gcc -O0)       & 390          & 391          & 396          & 416          & 390 & 911 \\
Python3           & 7{,}650      & 7{,}722      & 7{,}761      & 8{,}316      & 7{,}726 & 9{,}054 \\
\bottomrule
\end{tabular}
\end{table}

Aricode achieves execution times within 7\% of hand-optimized assembly across all benchmarks and outperforms GCC~-O3 by a factor of 1.2--2.6$\times$. The Mersenne challenge shows the largest gap with GCC~-O0 (353~$\mu$s vs.\ 911~$\mu$s), demonstrating the effectiveness of tail call optimization for deep recursion.

\subsection{Results: Overall Ranking}

\begin{table}[H]
\centering
\caption{Overall sparring ranking (points system)}
\label{tab:ranking}
\begin{tabular}{@{}clr@{}}
\toprule
\textbf{Rank} & \textbf{Implementation} & \textbf{Points} \\
\midrule
\textbf{1} & \textbf{Aricode} & \textbf{132} \\
2 & NASM x86\_64 assembly & 130 \\
3 & C (gcc -O2) & 74 \\
4 & C (gcc -O0) & 70 \\
5 & C (gcc -O3) & 67 \\
6 & C (gcc -Os) & 63 \\
7 & Python 3.12 & 56 \\
8 & C (gcc -O2 dynamic) & 56 \\
\bottomrule
\end{tabular}
\end{table}

Aricode achieves overall first place with 132 points, narrowly ahead of hand-written NASM assembly (130 points). The advantage stems from consistently smaller binary sizes while maintaining near-identical execution speeds. GCC-compiled C, despite decades of optimizer development, places third with roughly half the points, primarily penalized by large binary sizes from CRT and libc overhead.

\subsection{Impact of Tail Call Optimization}

The Mersenne prime verification challenge provides a clear demonstration of tail call optimization impact:

\begin{table}[H]
\centering
\caption{Mersenne M31 verification: effect of tail call optimization}
\label{tab:tco}
\begin{tabular}{@{}lrrr@{}}
\toprule
\textbf{Configuration} & \textbf{Exec Time ($\mu$s)} & \textbf{Stack Usage} & \textbf{Improvement} \\
\midrule
Aricode (no TCO)  & 836 & $O(n)$, 23{,}170 frames & --- \\
\textbf{Aricode (with TCO)} & \textbf{352} & $O(1)$, 1 frame & \textbf{57\%} \\
NASM (iterative)  & 310 & $O(1)$ & reference \\
C (gcc -O3, recursive) & 427 & $O(n)$, 23{,}170 frames & --- \\
\bottomrule
\end{tabular}
\end{table}

With TCO, Aricode's recursive implementation achieves performance within 14\% of NASM's hand-written iterative loop, while GCC~-O3 fails to optimize the recursive C implementation and performs 21\% slower than Aricode.

% ============================================================
\section{Related Work}
\label{sec:related}

\paragraph{Error Handling in Programming Languages.}
Go~\cite{go2012} returns error values as a second return parameter but does not enforce checking. Rust~\cite{rust2015} uses the \texttt{Result$\langle$T,E$\rangle$} type with the \texttt{?} operator, but allows \texttt{unwrap()} which can panic. Swift uses optionals with forced unwrapping. Aricode's approach is more restrictive: there is no mechanism to bypass error handling at compile time.

\paragraph{Minimal Binary Generation.}
Projects like \texttt{musl libc} and TinyCC~\cite{bellard2001} aim to reduce binary sizes. Aricode eliminates the C runtime entirely, generating standalone ELF binaries with no external dependencies.

\paragraph{Direct Code Generation.}
QBE~\cite{qbe} provides a lightweight backend alternative to LLVM. Aricode goes further by generating machine code directly from the AST without any intermediate representation.

\paragraph{Decimal Arithmetic.}
IBM's General Decimal Arithmetic specification~\cite{ieee754r} standardizes decimal floating-point. Python's \texttt{decimal} module and Java's \texttt{BigDecimal} provide library-level support. Aricode integrates decimal arithmetic as a native type with compile-time guarantees.

% ============================================================
\section{Limitations and Future Work}
\label{sec:future}

Current limitations include: the codegen targets only Linux x86\_64; the type system lacks user-defined types (structs, enums); loop constructs (\texttt{while}, \texttt{for}) are parsed but not yet compiled (recursion with TCO substitutes); and the standard library is minimal. Tail call optimization currently applies only to self-recursive calls, not mutual recursion.

Future work includes: ARM64 and WebAssembly backends; struct and enum types; a module system with package manager; SIMD vectorization for numeric workloads; and inter-procedural optimization.

% ============================================================
\section{Conclusion}
\label{sec:conclusion}

Aricode demonstrates that mandatory error handling and competitive performance are not mutually exclusive. By classifying all possible failures into five levels and enforcing handling at compile time, Aricode eliminates the entire category of silent runtime errors---the most dangerous class of software defects.

The direct x86\_64 code generator, combined with constant folding, peephole optimization, and tail call optimization, produces binaries that are 2{,}000$\times$ smaller than GCC output and execute 1.2$\times$ faster than GCC~-O3. In a comprehensive benchmark suite including Ackermann function evaluation, Collatz conjecture verification, and Mersenne prime testing, Aricode achieves first place overall against eight implementations including hand-optimized x86\_64 assembly.

The compiler, benchmark suite, and all source code are available at: \url{https://github.com/lynxcraft/aricode}.

% ============================================================
\begin{thebibliography}{99}

\bibitem{ariane5}
J.-M. Jazequel and B. Meyer, ``Design by Contract: The Lessons of Ariane,'' \textit{IEEE Computer}, vol.~30, no.~1, pp.~129--130, 1997.

\bibitem{toyota2013}
M. Barr, ``Bookout v. Toyota Motor Corp.,'' Expert Witness Report on Software Analysis, 2013.

\bibitem{collatz1937}
L. Collatz, ``On the motivation and origin of the $(3n+1)$ problem,'' \textit{Journal of Qufu Normal University}, Natural Science Edition, 1986 (original conjecture 1937).

\bibitem{go2012}
R. Pike, ``Go at Google: Language Design in the Service of Software Engineering,'' in \textit{Proc.\ 3rd Annual Conference on Systems, Programming, and Applications (SPLASH)}, 2012.

\bibitem{rust2015}
N. Matsakis and F. Klock, ``The Rust Language,'' in \textit{Proc.\ ACM SIGAda Ada Letters}, vol.~34, no.~3, pp.~103--104, 2014.

\bibitem{bellard2001}
F. Bellard, ``TCC: Tiny C Compiler,'' 2001. Available: \url{https://bellard.org/tcc/}

\bibitem{qbe}
Q. Carbonneaux, ``QBE: A Compiler Backend,'' 2017. Available: \url{https://c9x.me/compile/}

\bibitem{ieee754r}
IEEE, ``IEEE 754-2008: Standard for Floating-Point Arithmetic,'' 2008.

\end{thebibliography}

\end{document}
